\chapter{Methodology and Proposed Framework}
\label{ch:methodology}



Chapter~\ref{ch:introduction} framed the central question of this thesis: whether perception built solely on stereo cameras can be made accurate enough that sharing detections between two cooperating vehicles produces a measurable improvement in what each one perceives. This chapter answers the how. It presents the proposed system through two complementary views: a structural view, decomposing the perception problem into a four-stage pipeline in which each stage carries a single responsibility, a documented interface to its neighbours, and a standalone validation step; and a procedural view, tracing the V-model development process that keeps design, implementation, and verification in deliberate correspondence, so that confidence in the final cooperative result rests on stage-by-stage evidence rather than on end-to-end output alone.




\section{Design Requirements and Principles}
\label{sec:requirements}

Before any architecture can be proposed, the objectives of
Chapter~\ref{ch:introduction} must be translated into concrete engineering
requirements. These fall into functional requirements (what the system must do) and
non-functional requirements (the qualities it must exhibit).

\paragraph{Functional requirements.}
\begin{itemize}[leftmargin=*,itemsep=3pt]
  \item \textbf{FR1 --- Depth from stereo.} Given a rectified stereo image pair
  from one vehicle, the system shall produce a dense per-pixel estimate of scene
  depth.
  \item \textbf{FR2 --- Object detection.} The system shall detect the surrounding
  vehicles in a single camera image as two-dimensional regions.
  \item \textbf{FR3 --- Lifting to 3D.} The system shall combine detections with
  depth to place each detected vehicle at a three-dimensional position in the
  observing vehicle's coordinate frame.
  \item \textbf{FR4 --- Cooperative fusion.} Given the independent 3D observations
  of two vehicles and their relative pose, the system shall register both sets of
  observations into a common frame and combine them into a single, consistent
  scene.
  \item \textbf{FR5 --- Symmetric benefit measurement.} The system shall make it
  possible to quantify, for \emph{each} of the two vehicles, the benefit obtained
  from cooperation --- both the objects recovered that the vehicle could not see
  alone and any change in localization accuracy.
\end{itemize}

\paragraph{Non-functional requirements.}
\begin{itemize}[leftmargin=*,itemsep=3pt]
  \item \textbf{NFR1 --- Modularity.} Each perception stage shall be independently
  runnable, testable, and replaceable, communicating with its neighbours only
  through a documented data interface.
  \item \textbf{NFR2 --- Reproducibility.} Every run shall be deterministic given
  its configuration, and all parameters and outcomes shall be recorded so that any
  result can be regenerated and audited.
  \item \textbf{NFR3 --- Stagewise verifiability.} No stage shall depend on the
  correctness of a later stage for its own validation; each shall be verifiable
  against ground truth in isolation.
  \item \textbf{NFR4 --- Honesty of output.} The system shall report only those
  quantities that the sensing modality can reliably support, and shall explicitly
  decline to report those it cannot.
\end{itemize}

\section{The V-Model Development Process}
\label{sec:vmodel}

The system was developed under a V-model lifecycle. The V-model is well
suited to this work for two reasons. First, the perception pipeline is naturally
decomposable into stages whose interfaces can be specified up front, which favours
a process that fixes requirements and design before implementation. Second, and
more importantly, the thesis question is ultimately a question of trust: a
cooperative gain is only meaningful if each link in the chain that produces it has
been independently verified. The V-model makes this verification structural rather
than incidental --- every level of design is paired, from the outset, with the
level of testing that will confirm it.

The left arm of the V descends from abstract intent to concrete code: research
objectives are refined into the system requirements of Section~\ref{sec:requirements},
the requirements into the architecture of Section~\ref{sec:architecture}, the
architecture into the detailed design of each individual stage, and the stage
designs into implementation. The right arm ascends back up through progressively
broader levels of testing, and --- this is the defining feature of the model ---
each ascending verification level is tied horizontally to the descending design
level it confirms (Figure~\ref{fig:vmodel}).

\begin{figure}[htbp]
\centering
\hyphenpenalty=10000\exhyphenpenalty=10000
\resizebox{\textwidth}{!}{%
\begin{tikzpicture}[node distance=0pt,vbox/.append style={minimum width=3.6cm,text width=3.3cm}]
% Left (design) arm — descending
\node[vbox,fill=blue!6]   (L1) at (0,8)   {Requirements / Objectives\\\footnotesize(Ch.~1)};
\node[vbox,fill=blue!6]   (L2) at (1.6,6)  {System Architecture\\\footnotesize(Sec.~\ref{sec:architecture})};
\node[vbox,fill=blue!6]   (L3) at (3.2,4)  {Stage / Module Design\\\footnotesize(Sec.~\ref{sec:stage1}--\ref{sec:stage4})};
\node[vbox,fill=blue!6]   (L4) at (4.8,2)  {Detailed Design\\\footnotesize(stage modules)};
% Right (verification) arm — ascending
\node[vbox,fill=green!8]  (R1) at (14.4,8) {Acceptance Validation\\\footnotesize coop.\ dynamics \& limits};
\node[vbox,fill=green!8]  (R2) at (12.8,6) {System Testing\\\footnotesize end-to-end agent pair};
\node[vbox,fill=green!8]  (R3) at (11.2,4) {Stage Validation\\\footnotesize vs.\ ground truth (NFR3)};
\node[vbox,fill=green!8]  (R4) at (9.6,2)  {Unit Testing\\\footnotesize per-stage automated};
% Bottom vertex
\node[vbox,fill=gray!12]  (IM) at (7.2,0.2) {Implementation};
% V arms
\draw[flow] (L1) -- (L2); \draw[flow] (L2) -- (L3); \draw[flow] (L3) -- (L4); \draw[flow] (L4) -- (IM);
\draw[flow] (IM) -- (R4); \draw[flow] (R4) -- (R3); \draw[flow] (R3) -- (R2); \draw[flow] (R2) -- (R1);
% Horizontal correspondences
\draw[corr] (L1) -- (R1); \draw[corr] (L2) -- (R2); \draw[corr] (L3) -- (R3); \draw[corr] (L4) -- (R4);
\end{tikzpicture}}
\caption[The V-model as instantiated in this project]{The V-model as instantiated
in this project.}
\label{fig:vmodel}
\end{figure}

The horizontal correspondences are realized concretely in this project:

\begin{itemize}[leftmargin=*,itemsep=4pt]
  \item \textbf{Implementation $\leftrightarrow$ Unit testing.} Each stage is
  accompanied by an automated test module exercising its internal logic in isolation from the rest of the pipeline and from
  real data. These tests encode the detailed design contracts and are the first
  gate any change must pass.
  \item \textbf{Stage design $\leftrightarrow$ Stage validation.} Each stage has a
  dedicated, standalone validation procedure that scores its output against ground
  truth without invoking later stages (NFR3). This is the level at which a stage's
  design --- not just its code --- is confirmed to behave as intended on
  real inputs.
  \item \textbf{System architecture $\leftrightarrow$ System testing.} The
  architecture is exercised as a whole by running the complete chain --- depth,
  detection, lifting, and fusion --- for a genuine pair of cooperating agents,
  confirming that the stage interfaces compose correctly end to end.
  \item \textbf{Objectives $\leftrightarrow$ Acceptance validation.} The top-level
  question is answered by the final measurement of the cooperative gain in a
  controlled multi-agent setting, which is the engineering acceptance criterion for
  the thesis as a whole.
\end{itemize}

A practical consequence of this discipline is that defects are localized to the
narrowest possible level. A regression in a geometric routine surfaces at the unit
level; a degradation in depth quality surfaces at stage validation; an interface
mismatch surfaces at system testing. Because each level has an explicit owner in
the design, failures are diagnosed against the corresponding specification rather
than against the opaque end-to-end behaviour.

\section{System Architecture Overview}
\label{sec:architecture}

The proposed system is a \textbf{four-stage pipeline}. Each stage consumes the
output of its predecessor through a documented interface and emits artifacts that
are both the input to the next stage and an independently inspectable result. The
stages and their data interfaces are summarized in Table~\ref{tab:stages}, and their
end-to-end composition across a pair of cooperating vehicles is shown in
Figure~\ref{fig:pipeline}.

\begin{table}[htbp]
\centering
\caption{The four pipeline stages and their interfaces.}
\label{tab:stages}
\begin{tabularx}{\textwidth}{@{}clL@{}}
\toprule
\textbf{Stage} & \textbf{Name} & \textbf{Input $\rightarrow$ Output} \\
\midrule
1 & Depth  & Rectified stereo pair $\rightarrow$ dense disparity / metric depth \\
2 & Detect & Single (left) image $\rightarrow$ 2D vehicle regions \\
3 & Lift   & 2D regions + depth $\rightarrow$ 3D position per vehicle  \\
4 & Fuse & Two agents' 3D observations + relative pose $\rightarrow$ one fused, registered scene \\
\bottomrule
\end{tabularx}
\end{table}

Stages 1 through 3 constitute the \textbf{single-vehicle perception chain}: they
take one vehicle's raw stereo imagery and turn it into a set of 3D vehicle
positions expressed in that vehicle's own camera frame. Stage 4 is the
\textbf{cooperative layer}: it takes the perception output of two such vehicles and
the geometric relationship between them, and produces a single shared scene. The first three stages fulfill Objective~1 (the single-agent 3D localisation baseline); Stage 4 addresses Objective~2 (the decentralised object-level fusion layer); and the performance analysis of the combined, end-to-end system addresses Objective~3 (the empirical characterisation of cooperative localisation dynamics).


\begin{figure}[htbp]
\centering
\resizebox{\textwidth}{!}{%
\begin{tikzpicture}[
  databox/.append style={minimum width=2.6cm},
  stagebox/.append style={minimum width=2.9cm},
  lbl/.style={font=\scriptsize,fill=gray!12,inner sep=1.5pt,rounded corners=1pt}]
% ---- Agent A : S1 (depth) and S2 (detection) run in parallel ----
\node[databox]  (srcA)  at (0,6)    {Stereo pair (L,R)};
\node[databox]  (imgA)  at (3.6,6)  {Left image};
\node[stagebox] (s1A)   at (0,4.3)  {S1 — Depth\\\footnotesize disparity $\rightarrow$ depth};
\node[stagebox] (s2A)   at (3.6,4.3){S2 — Detection\\\footnotesize 2D vehicle boxes};
\node[stagebox] (s3A)   at (1.8,2.6){S3 — Lift to 3D\\\footnotesize 3D positions};
\node[font=\bfseries\small] at (1.8,7) {Vehicle A};
% ---- Agent B (mirror) ----
\node[databox]  (srcB)  at (9,6)    {Stereo pair (L,R)};
\node[databox]  (imgB)  at (12.6,6) {Left image};
\node[stagebox] (s1B)   at (9,4.3)  {S1 — Depth\\\footnotesize disparity $\rightarrow$ depth};
\node[stagebox] (s2B)   at (12.6,4.3){S2 — Detection\\\footnotesize 2D vehicle boxes};
\node[stagebox] (s3B)   at (10.8,2.6){S3 — Lift to 3D\\\footnotesize 3D positions};
\node[font=\bfseries\small] at (10.8,7) {Vehicle B};
% ---- intra-agent flows ----
\foreach \s/\d in {srcA/s1A,imgA/s2A,srcB/s1B,imgB/s2B} \draw[flow] (\s) -- (\d);
\draw[flow] (s1A.south) -- (s3A.north) ;
\draw[flow] (s2A.south) -- (s3A.north) ;
\draw[flow] (s1B.south) -- (s3B.north) ;
\draw[flow] (s2B.south) -- (s3B.north) ;
% ---- fusion ----
\node[fusebox] (fus) at (6.3,0.7)
  {S4 — V2V Fusion\\\footnotesize register $\cdot$ match $\cdot$ merge};
\draw[flow] (s3A.south) |- (fus.west);
\draw[flow] (s3B.south) |- (fus.east);
\node[databox,fill=orange!6] (pose) at (10.0,0.7) {relative pose\\\footnotesize $T_{B\rightarrow A}$};
\draw[flow] (pose) -- (fus);
\node[databox,fill=orange!6,minimum width=5cm] (out) at (6.3,-1.4)
  {Fused scene in A's frame\\\footnotesize (symmetric gain measurable)};
\draw[flow] (fus) -- (out);
\end{tikzpicture}}
\caption[End-to-end data flow]{End-to-end data flow of the four-stage pipeline.}
\label{fig:pipeline}
\end{figure}

Two architectural properties follow from the requirements of
Section~\ref{sec:requirements} :

\begin{itemize}[leftmargin=*,itemsep=3pt]
  \item \textbf{Symmetry of the perception chain.} The identical Stage 1--3 chain
  is applied to each vehicle. Neither agent is privileged; cooperation is a
  relationship between two equal perception systems, which is what makes a
  symmetric benefit measurement (FR5) meaningful.
  \item \textbf{A single point of coupling.} The two perception chains are entirely
  independent until Stage 4. All inter-agent coupling is confined to the
  fusion layer, keeping the single-vehicle design free of cooperation-specific
  concerns.
\end{itemize}

\section{Data Sources and the Dual-Dataset Strategy}
\label{sec:data}

A recurring difficulty identified in Chapter~\ref{ch:introduction} is that no
single dataset offers both real stereo imagery and simultaneous
multi-vehicle recordings with trustworthy ground truth. The architecture resolves
this with a deliberate dual-dataset strategy, assigning each data source
the role it is uniquely suited to.

\begin{itemize}[leftmargin=*,itemsep=4pt]
  \item \textbf{Real-image stereo source (single-vehicle chain).} Stages 1--3 are
  designed and validated against real-world driving stereo data. This is the only
  way to confirm that the camera-only perception chain behaves on genuine imagery,
  with its real noise, lighting, and scene complexity. Because this source provides
  one vehicle's view at a time, it cannot exercise cooperation, but it is the
  authoritative ground for the single-vehicle claim.
  \item \textbf{Simultaneous multi-agent source (cooperative layer).} Stage 4
  requires two vehicles observing the same scene at the same instant,
  together with exact relative pose and a complete inventory of which objects each
  vehicle can truly see. These conditions are obtainable from a controlled
  simulation environment~\cite{dosovitskiy2017carla} that records two moving ego
  vehicles in a shared scene, with per-object, per-agent visibility derived from
  the rendered view rather than guessed geometrically. This source is what makes
  the cooperative gain measurable (FR5) with reliable ground truth.
\end{itemize}

This division is an engineering decision rather than a compromise: each stage is
exercised on the data that can actually test it. The shared
interface contract between the stages (Table~\ref{tab:stages}) is what allows the
same Stage 1--3 design to operate on either source.

\section{Stage 1 --- Stereo Depth Estimation}
\label{sec:stage1}

The remaining sections of this chapter make use of mathematical notation throughout. To keep them self-contained, the notation used is collected once in Appendix~\ref{app:notation} (Table~\ref{tab:notation}).

\paragraph{Purpose.} Stage 1 converts a rectified stereo image pair into a dense
estimate of per-pixel disparity, from which metric depth is recovered by the
standard stereo triangulation relationship~\cite{szeliski2022computer}
\begin{equation}
Z = \frac{f \cdot B}{d}
\label{eq:triangulation}
\end{equation}
where $Z$ is the depth of a pixel, $f$ the focal length in pixels, $B$ the stereo
baseline in metres, and $d$ the disparity (the horizontal displacement of the pixel
between the left and right images). The baseline itself is recovered from the two
cameras' projection matrices, $B = (t_x^{L} - t_x^{R}) / f$, where $t_x^{L}$ and
$t_x^{R}$ are the horizontal translation terms of the left and right projection
matrices.

This relationship is central to the entire methodology and reappears as the
dominant source of error analysed in later chapters. Differentiating it with
respect to disparity gives the depth error induced by a disparity error:
\begin{equation}
\left| \Delta Z \right| \;\approx\; \frac{Z^{2}}{f \cdot B}\, \left| \Delta d \right|
\label{eq:deptherror}
\end{equation}

Because the depth error scales with the square of the distance, a fixed
disparity error of a fraction of a pixel translates into a small error nearby but a
large one far away. The system is therefore expected, by construction, to be most
accurate on nearby vehicles --- a property that informs the design of every
subsequent stage and that is quantified empirically in Chapter~5.

\paragraph{Design.} The stage is built around two interchangeable depth estimators
that represent two distinct engineering trade-offs, selectable without altering any
downstream stage:

\begin{itemize}[leftmargin=*,itemsep=3pt]
  \item A \textbf{classical semi-global block-matching estimator}~\cite{hirschmuller2008},
  which produces a sparse disparity map: it reports depth only where a
  confident left--right correspondence exists and leaves textureless or ambiguous
  regions unfilled. Its sparsity is, in effect, a built-in reliability filter ---
  the pixels it does report tend to lie on well-defined surfaces.
  \item A \textbf{learned deep-stereo estimator}~\cite{wang2026}, which
  produces a dense disparity map with full coverage, at greater
  computational cost. It is run directly as an in-process inference step.
\end{itemize}

Offering both is a deliberate architectural choice. It allows the methodology to
separate questions of depth accuracy from questions of downstream
robustness: a sparse, conservative map and a dense, complete map stress the later
stages differently, and carrying both through the pipeline lets the project reason
about that trade-off rather than presuppose a single answer.

\paragraph{Interface.} The stage emits a per-pixel disparity map in which invalid
pixels are explicitly marked, so that downstream stages can distinguish ``no
measurement'' from ``a measurement of zero.'' This explicit invalid-pixel contract
is what allows the sparse and dense estimators to share one interface.

\section{Stage 2 --- 2D Object Detection}
\label{sec:stage2}

\paragraph{Purpose.} Stage 2 locates the surrounding vehicles in the (left) camera
image as two-dimensional regions. It answers where in the image the objects
of interest are, the question of where in space is addressed separately, in Stage 3.

\paragraph{Design.} The stage uses a pretrained transformer-based object
detector~\cite{lv2024} applied to the left image of the stereo pair. Detection
operates on the single image rather than on the depth map because appearance-based
detection is mature and robust, whereas the depth map is better used for
localizing an already-detected object than for finding it. The
detector is pretrained on a general-purpose object vocabulary; the stage maps the
relevant vehicle categories of that vocabulary onto the project's target class and
discards the rest, applying a confidence threshold to suppress weak detections.

\paragraph{Class scope.} The pipeline is deliberately scoped to a single object
class --- the car. Restricting the class keeps the evaluation interpretable and aligned
with the honesty principle (NFR4). The architecture does not preclude additional
classes; re-enabling them is a configuration concern localized to the detection and
matching stages.

\paragraph{Interface.} The stage emits, per image, a set of labelled 2D regions
with associated confidences, which become the input to the lifting stage alongside
the Stage 1 depth map.

\section{Stage 3 --- Lifting to 3D Position}
\label{sec:stage3}

\paragraph{Purpose.} Stage 3 is the junction of the two preceding stages: it
combines each 2D detection (Stage 2) with the depth map (Stage 1) and the camera
calibration to place the detected vehicle at a three-dimensional position in the
observing vehicle's camera frame. This is the stage at which the system crosses
from image space into the metric 3D world in which cooperation takes place.

\paragraph{The position-only design decision.} The single most important
methodological decision in the single-vehicle chain is that Stage 3 outputs 3D
position only --- a point $(x, y, z)$ per vehicle, carried together with the source
2D region --- and deliberately does not output a full 3D bounding box with
size and heading. This decision is a direct application of the honesty principle
(NFR4) and follows from the range-dependent stereo uncertainty at the heart of the
problem statement (Chapter~\ref{ch:introduction}).

The reasoning is geometric. Recovering an object's position requires
establishing how far away it is, which stereo supports. Recovering an object's
size and heading, by contrast, requires resolving its
three-dimensional extent and orientation from a single viewpoint at range --- and
the same disparity-to-depth sensitivity that limits position accuracy at distance
makes extent and orientation far less observable still. 

\paragraph{Design.} For each 2D region, the stage samples a representative depth
from the disparity map within the region, converts it to a metric distance via the
stereo relationship of Section~\ref{sec:stage1}, and unprojects the region's centre
to a 3D point using the camera calibration. All 3D quantities are expressed in the
KITTI camera convention ($X$ right, $Y$ down, $Z$ forward into the scene;
Figure~\ref{fig:frames}). Given a region centre at pixel $(u, v)$ and its sampled
depth $Z$, the inverse of the pinhole projection yields the 3D point
\begin{equation}
X = \frac{(u - c_x)\, Z}{f_x}, \qquad
Y = \frac{(v - c_y)\, Z}{f_y}, \qquad
Z = Z
\label{eq:unproject}
\end{equation}
where $(f_x, f_y)$ are the focal lengths and $(c_x, c_y)$ the principal point of the
left camera, both read from its calibration. The point $(X, Y, Z)$ is the vehicle's
position in the observing vehicle's camera frame. Two considerations shape this:

\begin{itemize}[leftmargin=*,itemsep=4pt]
  \item \textbf{Robust depth sampling within a region.} A detection's region
  contains pixels belonging to the vehicle but also, typically, road surface and
  background. The stage therefore does not take a naive average; it samples depth in
  a way designed to favour the object surface over contaminating background, with
  the sampling strategy tuned per depth estimator because the sparse and dense maps
  of Stage 1 have very different pixel distributions inside a region. The specific
  tuning is treated as an experimental matter and reported in Chapter~5.
  \item \textbf{Confidence propagation.} The 3D position inherits a confidence that
  scales the 2D detection confidence by the proportion of the region for which valid
  depth was available, $c_{3D} = c_{2D} \cdot \rho$, where the coverage ratio
  $\rho \in [0, 1]$ is the fraction of region pixels carrying a valid depth
  measurement. A position lifted from a well-covered region is thus trusted more
  than one lifted from a sparsely measured region, and this confidence later
  participates directly in the fusion merge (Section~\ref{sec:stage4}).
\end{itemize}

\paragraph{Interface.} The stage emits, per detected vehicle, a labelled 3D
position with its propagated confidence and the source 2D region. This compact,
position-only representation is intentionally schema-compatible with the fusion
core, which is designed to operate whether or not size and heading are present.

\section{Stage 4 --- Cooperative V2V Fusion}
\label{sec:stage4}

Stage 4 is the cooperative layer and the realization of the thesis goal. It takes
the independent 3D observations of two vehicles and combines them into a single
shared scene in which objects visible to either vehicle become known to both.
Combining detections after each vehicle has perceived independently is the
late fusion strategy for cooperative perception~\cite{xu2022opv2v}, chosen
here because it keeps the inter-vehicle interface compact (3D positions rather than
raw imagery) and the single-vehicle chain unaware of cooperation. Its design
separates a source-agnostic fusion core from the data plumbing that
feeds it, so that the same fusion logic operates on either real or simulated inputs
and on either position-only or fully-specified boxes.

The fusion proceeds in three conceptual steps --- \textbf{register, match, merge}
--- preceded by the construction of the common frame (Figure~\ref{fig:fusion}).

\paragraph{Common frame and registration.} Each vehicle perceives in its own camera
frame; the two frames are unrelated until the vehicles' relative pose is known.
Cooperation therefore begins by choosing one vehicle's frame (Vehicle A's) as the
reference and transforming the other vehicle's observations into it by a single
rigid $4\times4$ homogeneous transform derived from the two vehicles' poses,
\begin{equation}
T_{B \rightarrow A} = M \; T_{W,A}^{-1} \; T_{W,B} \; M^{-1}
\label{eq:registration}
\end{equation}
where $T_{W,A}$ and $T_{W,B}$ are the world-from-camera transforms of the two
vehicles (each built from the vehicle's world pose composed with its fixed camera
mount), and $M$ is the constant axis map from the vehicle frame to the KITTI camera
convention, $(x, y, z)_{\text{cam}} = (y, -z, x)_{\text{vehicle}}$
(Figure~\ref{fig:frames}). Any 3D point observed by Vehicle B is then carried into
A's frame by $\tilde{p}_A = T_{B \rightarrow A}\, \tilde{p}_B$ (in homogeneous
coordinates), and a box's heading is offset by the transform's rotation about the
camera vertical axis, $\theta_{\text{offset}} = \operatorname{atan2}(R_{13}, R_{33})$,
extracted from the rotation block $R$ of $T_{B \rightarrow A}$. Because the
cooperative source records simultaneous observations, this is a purely spatial
alignment --- there is no temporal extrapolation, and a static object seen by both
vehicles should, after registration, coincide. The design exploits exactly this
property as a built-in correctness check: an object seen by both agents must
register to near-zero displacement, and a large displacement is treated as a symptom
of a coordinate or pose error rather than of a moving object. This sanity check is
part of the methodology: it is the mechanism by which the
most dangerous class of bug in cooperative perception, a silent frame misalignment,
is caught.

\begin{figure}[htbp]
\centering
\begin{tikzpicture}[
    >={Latex[length=2.4mm]},
    axis/.style={->,thick,black!82},
    ttl/.style={font=\bfseries\small,align=center},
    sub/.style={font=\footnotesize,text=black!55,align=center},
    alab/.style={font=\small},
    panel/.style={draw=black,rounded corners=4pt,minimum width=5.6cm,minimum height=5.7cm}]
% ===== equal-size frames (black border, no fill) =====
\begin{scope}[on background layer]
  \node[panel] (cpan) at (1.4,0.45) {};
  \node[panel] (vpan) at (10.0,0.45) {};
\end{scope}
% ===== KITTI camera frame axes =====
\coordinate (oc) at (0,0);
\draw[axis] (oc) -- (2.0,0)    node[alab,right]      {$X$ (right)};
\draw[axis] (oc) -- (0,-1.7)   node[alab,below]      {$Y$ (down)};
\draw[axis] (oc) -- (1.2,0.95) node[alab,above right]{$Z$ (forward)};
\fill[black!82] (oc) circle (2.2pt);
\node[ttl] at (1.4,2.7)  {KITTI camera frame};
\node[sub] at (1.4,2.28) {optical, right-handed};
% ===== Vehicle frame axes =====
\coordinate (ov) at (8.9,0);
\draw[axis] (ov) -- ($(ov)+(1.2,0.95)$) node[alab,above right]{$x$ (forward)};
\draw[axis] (ov) -- ($(ov)+(2.0,0)$)    node[alab,right]      {$y$ (right)};
\draw[axis] (ov) -- ($(ov)+(0,1.6)$)    node[alab,left]       {$z$ (up)};
\fill[black!82] (ov) circle (2.2pt);
\node[ttl] at (10.0,2.7)  {Vehicle (agent) frame};
\node[sub] at (10.0,2.28) {left-handed, CARLA};
% ===== axis-map bridge arrow, in the gap between the two boxes =====
\coordinate (mid) at (0,0.5);
\draw[-{Latex[length=2.6mm]},thick,black!55]
  (vpan.west |- mid) to[bend right=14]
  node[above=1pt,font=\footnotesize\itshape,text=black!60]{axis map $M$}
  (cpan.east |- mid);
% ===== axis-map equation, centred under the whole figure =====
\node[draw=black!55,rounded corners=3pt,fill=none,font=\footnotesize,inner sep=7pt,anchor=north]
  at ([yshift=-0.45cm]current bounding box.south)
  {Axis map\quad $(X,\,Y,\,Z)_{\text{cam}} = (\,y,\; -z,\; x\,)_{\text{vehicle}}$};
\end{tikzpicture}
\caption[The two coordinate frames]{Spatial Mapping Between KITTI Optical and CARLA Vehicle Coordinate Frames}
\label{fig:frames}
\end{figure}

\paragraph{Matching.} Registered observations from the two vehicles are associated
greedily by bird's-eye-view (ground-plane) centre distance, within object class,
accepting a pair as the same physical object only when their separation falls below
a per-class threshold $\tau$. The bird's-eye-view distance discards the vertical
axis and measures separation in the ground plane,
\begin{equation}
d_{\text{BEV}}(a, b) = \sqrt{(x_a - x_b)^2 + (z_a - z_b)^2}, \qquad
\text{match if } d_{\text{BEV}} \le \tau
\label{eq:bev}
\end{equation}
candidate pairs being considered in ascending order of distance so that the closest
compatible pair is committed first. Matching on ground-plane centre distance ---
rather than on 3D box overlap --- is the natural choice given the position-only
output of Stage 3, and it is exactly what makes the fusion core able to consume both
position-only and fully-specified inputs through one code path.

\paragraph{Merge.} A corroborated pair --- one object that both vehicles saw --- is
merged into a single estimate. Writing the two observations' confidences as $c_A$
and $c_B$, each merged position coordinate is the confidence-weighted average of the
two observations, and the fused confidence is combined by a noisy-OR rule:
\begin{equation}
\hat{x} = \frac{c_A\, x_A + c_B\, x_B}{c_A + c_B}
\quad (\text{and likewise for } y, z), \qquad
\hat{c} = 1 - (1 - c_A)(1 - c_B)
\label{eq:merge}
\end{equation}
The noisy-OR rule reflects that two independent detections of the same object are
more trustworthy than either alone (the fused confidence exceeds both inputs). Where
the inputs happen to carry size and heading, those are merged too --- dimensions by
the same confidence-weighted average, and heading by a confidence-weighted
\emph{circular} mean,
\begin{equation}
\hat{\theta} = \operatorname{atan2}\!\Big(\textstyle\sum_i c_i \sin\theta_i,\; \sum_i c_i \cos\theta_i\Big),
\label{eq:circmean}
\end{equation}
which averages angles correctly across the $\pm\pi$ wrap-around. The merge is
nonetheless designed to degrade gracefully to position-only when size and heading
are absent. An object seen by only one vehicle is carried into the shared scene
unmerged, tagged with its source --- this is precisely the blind-spot recovery that
motivates the whole system. A matched pair whose post-registration displacement is
implausibly large is not merged but flagged, on the principle that for
simultaneous observations a large displacement indicates a faulty match or pose
error rather than genuine motion.

\begin{figure}[H]
\centering
\begin{tikzpicture}[node distance=0.6cm and 1.2cm]
\node[databox,fill=blue!8] (a) {A's observations\\\footnotesize $a_1\ a_2\ a_3$};
\node[databox,fill=green!8,right=3.2cm of a] (b)
  {B's observations\\\footnotesize (registered into A's frame)};
\node[fusebox,below=1.2cm of {$(a)!0.5!(b)$},minimum width=4.4cm] (match)
  {Match by BEV centre distance\\\footnotesize within class, thresholded $\le \tau$};
\draw[flow] (a.south) |- (match.west);
\draw[flow] (b.south) |- (match.east);
% outcomes
\node[databox,below left=1.3cm and -0.3cm of match,fill=orange!10,text width=3.0cm] (corr)
  {corroborated pair\\\footnotesize merge position,\\ noisy-OR confidence};
\node[databox,below=1.3cm of match,fill=gray!8,text width=2.8cm] (onlyA)
  {only-A object\\\footnotesize (recovered for B)};
\node[databox,below right=1.3cm and -0.3cm of match,fill=gray!8,text width=2.8cm] (onlyB)
  {only-B object\\\footnotesize (recovered for A)};
\foreach \n in {corr,onlyA,onlyB} \draw[flow] (match) -- (\n);
\node[fusebox,below=4.6cm of match,minimum width=6.2cm,fill=orange!12] (scene)
  {SHARED FUSED SCENE (A's frame)\\\footnotesize corroborated objects $+$ objects either agent saw alone};
\foreach \n in {corr,onlyA,onlyB} \draw[flow] (\n) -- (scene);
\end{tikzpicture}
\caption[The register--match--merge logic of the fusion core]{The
register--match--merge logic of the fusion core.}
\label{fig:fusion}
\end{figure}

\paragraph{Symmetric benefit by design.} The fused output is a single shared
scene, and this is what allows the cooperative benefit to be measured symmetrically
(FR5). Evaluation is framed against a cooperative ground truth --- the union
of every object visible to either vehicle, de-duplicated by object identity --- and
the shared scene is scored from both vehicles' standpoints: how many objects each
vehicle recovered that it could not have seen alone, and whether its localization of
objects it could see improved. The architecture treats the two vehicles identically
in this accounting, so neither agent's gain is privileged. Because the two ego
vehicles are themselves part of the simulated scene yet are not perception
targets (their poses are shared directly over the cooperative link), the
evaluation design treats an agent's detection of the other ego as neither a
hit nor a false alarm --- an ignore region --- so that the cooperation metric
reflects perceived traffic rather than the agents perceiving each other.

\bigskip
\noindent Two design commitments give the framework its character. The first is honesty
about sensing limits: because stereo cannot reliably recover object size or heading
at range, the system outputs 3D position only, and structures both its
representation and its evaluation around what cameras can actually provide. The
second is verification by construction: the V-model pairs every level of
design with the level of testing that confirms it, so that the headline cooperative
result rests on independently validated stages rather than on end-to-end behaviour
alone. With the architecture and process established, Chapter~4 details the
experimental setup and simulation environment, and Chapter~5 reports the stagewise
validation and the characterisation of the cooperative localisation dynamics that
this framework was built to measure.
